\section{Parametric Equations}
Parametric equations are when you are given multiple equations which use shared variables, and are calculated from the shared variables.

Commonly in exams this will be given in a form where you have $x$ and $y$ in terms of $t$ separately.

\begin{align*}
    x &= f(t)\\
    y &= g(t)
\end{align*}

You will normally be asked to find $y$ in terms of $x$. 
From here, there are two main approaches:
\begin{enumerate}
    \item Rearrange one of the equations for $t$ and substitute this into the other equation to eliminate $t$. This is sometimes not possible!
    \item Rearrange one of the equations so that it is in terms of the other equation, and then substitute the variable $x$ or $y$ in directly.
\end{enumerate}

We will see both approaches below.

\begin{example}
    A curve has equations:
    \begin{align*}
        x &= 2t + 1 \\
        y &= 4t - 5 + \frac{3}{t}, \; t\neq 0
    \end{align*}
    Let $x \neq 1$. Then show that the curve described by the above equations can be written as:

    \[ y = \frac{ax^2 + bx + c}{x-1}\]

    Where $a$, $b$, and $c$ are constants to be found.
\end{example}

\begin{solution}
    We can approach this by writing the simplest equation in terms of $t$. Note that we can't easily write $y$ in terms of $t$, so if we were not able to do this we would need to find another approach.

    \begin{align*}
        x &= 2t+1 \\
        x - 1 &= 2t \\
        \frac{x-1}{2} &= t
    \end{align*}

    And then substituting what $t$ is into $y$.

    \begin{align*}
        y &= 4t - 5 + \frac{3}{t} \\
         &= 4\frac{x-1}{2} - 5 +\frac{3}{(\frac{x-1}{2})} \tag{replacing each $t$ with $\frac{x-1}{2}$}\\
         &= 2(x-1) - 5 +\frac{3}{\frac{1}{2}(x-1)} \tag{taking out a factor of $\frac{1}{2}$} \\
         &= 2(x-1) - 5 + \frac{2}{2} \cdot \frac{3}{\frac{1}{2}(x-1)} \tag{multiplying by $1$} \\
         &= 2(x-1) - 5 + \frac{6}{\cancel{\frac{2}{2}}(x-1)} \tag{multiplying fractions}\\
         &= 2(x-1) - 5 + \frac{6}{x-1}\\
         &= \frac{2(x-1)(x-1)}{(x-1)} - \frac{5(x-1)}{(x-1)} + \frac{6}{x-1} \tag{common denominator}\\ 
         &= \frac{2(x-1)(x-1) - 5(x-1) + 6}{(x-1)} \tag{putting over one denominator}\\
         &= \frac{2(x^2-2x+1) - 5x+5 + 6}{(x-1)} \tag{expanding brackets}\\
         &= \frac{2x^2-4x+2 - 5x+11}{(x-1)} \tag{expanding brackets}\\
         &= \frac{2x^2-9x+13}{(x-1)} \tag{simplifying}
    \end{align*}
    Which finishes the question, finding $a=2$, $b=-9$, and $c=13$ so that 
    \[y = \frac{ax^2 + bx + c}{x-1} = \frac{2x^2+(-9)x+13}{(x-1)} \]

    \qed
\end{solution}

\begin{example}
    A curve has equations:
    \begin{align*}
        x &= \frac{1}{t} +5t\\
        y &= \frac{1}{t^2} + 10 + 25t^2, \; t\neq 0
    \end{align*}
    Show that the curve described by the above equations can be written as:

    \[ y = f(x)\]

    Where $f(x)$ is a function of $x$ to be found.
\end{example}

\begin{solution}
    We might want to approach this by writing the simplest equation in terms of $t$, but in this case it is not easily possible.
    For example, trying to rearrange the $x$ equation for $t$ you might get stuck at the point

    \[x = t(\frac{1}{t^2} +5)\]
    And not able to continue so simply.

    Instead, we attempt to rearrange the more complicated equation to be written in terms of the simpler one.

    What can we do with the second equation? Firstly, we can turn it into something that is easier to work with:

    \begin{align*}
        y &= \frac{1}{t^2} + 10 + 25t^2 \\
        \therefore t^2y &= 1 + 10t^2 + 25t^4
    \end{align*}
    Here, this almost looks familiar as a quadratic, except we have this pesky $t^2$ hanging around. Let's briefly use 
    a variable $u = t^2$ for convenience. Then 
    \begin{align*}
        t^2y &= 1 + 10t^2 + 25t^4 \\
        \therefore uy &= 1 + 10u + 25u^2 \tag{using $u = t^2$} \\
         &= (1 + 5u)(1 + 5u) \tag{factorising} \\
         &= (1+5u)^2\\
        \therefore uy &= (1+5u)^2 \tag{reminding ourselves} \\
        \therefore y &= \frac{1}{u}(1+5u)^2 \tag{rearranging}\\
        \therefore y &=  \frac{1}{t^2}(1+5t^2)^2 \tag{using $u = t^2$}\\
         &= (\frac{1}{t})^2(1+5t^2)^2 \tag{using $a^2/b^2 = (a/b)^2$}\\ 
         &= (\frac{1}{t}(1+5t^2))^2 \tag{using $a^2b^2 = (ab)^2$}\\ 
         &= (\frac{1}{t} + 5t)^2 \tag{expanding the bracket inside}
    \end{align*}

    This is as much simplification as we can get, and we have found that 
    \[ y = \frac{1}{t^2} + 10 + 25t^2  = (\frac{1}{t} + 5t)^2 \]

    Note that we could have gone straight to the answer if we noticed that we could factorise directly.

    But the last equality contains the equation for $x$ directly:

    \[x = \frac{1}{t} +5t\]

    Therefore we can replace any instance of $\frac{1}{t} +5t$ with $x$ to find:

    \[y = x^2\]

    This completes the question, finding $f(x) = x^2$.

    \qed
\end{solution}

\begin{example}
    A curve has equations:
    \begin{align*}
        x &= \cos(t)\\
        y &= \tan(t) + \sec(t)
    \end{align*}

    Where $0\leq t < \frac{\pi}{2}$.

    Show that the curve described by the above equations can be written as:

    \[ y = f(x) \]

    Where $f(x)$ is a function of $x$ to be found.

    Comment on whether the final function is the same as the curve for all $x$.
\end{example}

\begin{solution}
    We might want to approach this by writing the simplest equation in terms of $t$, but in this case we would need to invert $\cos$ and be careful with when this is allowed.

    Instead, we attempt to rearrange the more complicated equation to be written in terms of the simpler one.

    What can we do with the second equation? Firstly, we can write the definitions and then try to extract $\cos(t)$ wherever possible:

    \begin{align*}
        y &= \tan(t) + \sec(t)\\
         &= \frac{\sin t}{\cos t} + \frac{1}{\cos t}
    \end{align*}
    This is almost in terms of $\cos t$, except we have this pesky $\sin t$ lingering. We can make use of familiar equations to fix this (using that $0\leq t < \frac{\pi}{2}$):

    \begin{align*}
        \sin^2 t + \cos^2 t&= 1 \\
         \therefore \sin^2 t &= 1 - \cos^2 t \\
         \therefore \sin t &= \sqrt{1 - \cos^2 t}
    \end{align*}

    We can now substitute everything in to find:

    \begin{align*}
        y &= \frac{\sin t}{\cos t} + \frac{1}{\cos t} \\
         &= \frac{\sqrt{1 - \cos^2 t}}{\cos t } + \frac{1}{\cos t} \tag{substituting $\sin t$ above} \\
         &= \frac{\sqrt{1 - x^2}}{x} + \frac{1}{x} \tag{using $x = \cos t$}
    \end{align*}
    This finishes the question, finding 

    \[ y = f(x) = \frac{\sqrt{1 - x^2}}{x} + \frac{1}{x} \]

    Note that the question was careful to restrict $t$ between $0\leq t < \frac{\pi}{2}$. This means that only a portion of the \emph{full} curve given by 

    \[y  = \frac{\sqrt{1 - x^2}}{x} + \frac{1}{x}  \]
    is covered by the original equations for $x$ and $y$, specifically between the ranges of $x$ allowed by $\cos t$ inside this range, which is $0 < x \leq 1$.

    Therefore, the curve is actually the portion of the graph 

    \[y  = \frac{\sqrt{1 - x^2}}{x} + \frac{1}{x}  \]

    Where $0 < x \leq 1$. The expression defines a larger graph, but the parametric curve is only the part in this range of $x$.
    \qed
\end{solution}

The next question makes use of small angle approximations, covered in Section \ref{sec:trig-small-angle}.

\begin{example}
    Let $t>0$ be small, and measured in radians. A curve is given parametrically via 
    \begin{align*}
        x &= \cos(t^2) + 1 \\ 
        y &= \frac{1}{t^4} + t^8
    \end{align*}
    Using small angle approximations, show that the curve described by the above equations can be written as:

    \[ y \approx f(x) \]
    Where $f(x)$ is a function of $x$ to be found. Comment on what portion of the curve is valid for this approximation.
\end{example}

\begin{solution}
    There is no nice way to rearrange $y$ for $t$, and the question instructs us to use small angle approximations. Let's begin with recalling that if $\theta$ is small,

    \[\cos(\theta) \approx 1 - \frac{\theta^2}{2}\]

    And since $t$ is given to be small, we know that $t^2$ must be also small. Therefore,

    \[\cos(t^2) \approx 1 - \frac{(t^2)^2}{2}\]

    And therefore, we have that 

    \[ x = \cos(t^2) + 1 \approx 1 - \frac{t^4}{2} + 1  = 2 - \frac{t^4}{2}\]
    
    This is simple enough to rearrange for $t$:

    \begin{align*}
        x +   \frac{t^4}{2} &\approx 2 \\
         \therefore \frac{t^4}{2} &\approx 2 - x \\
         \therefore t^4 &\approx 4 - 2x \\
         \therefore t &\approx (4 - 2x)^{\frac{1}{4}}
    \end{align*}

    We can plug this into the equation for $y$:

    \begin{align*}
        y &= \frac{1}{t^4} + t^8 \\
         &\approx \frac{1}{((4 - 2x)^{\frac{1}{4}})^4} + ((4 - 2x)^{\frac{1}{4}})^8 \\
         &= \frac{1}{4 - 2x} + (4 - 2x)^2 \tag{using ${(a^b)}^c = a^{bc}$}\\
         &=\frac{1}{4 - 2x} + \frac{(4 - 2x)^2(4 - 2x)}{4 - 2x} \tag{common denominator}\\
         &=  \frac{1 + (4 - 2x)^3}{4 - 2x}
    \end{align*}

    This completes the question, as long as $x \neq 2$, which is always true. This is because the approximation is valid whenever $t$ is positive and small, which means that $x$ is close to $2$ but not equal to $2$ since $x = \cos(t^2) +1$ is only equal to $2$ at $t = 0$.
    Moreover, we deduce that $x$ must be slightly less than $2$ since $\cos(t^2)+1$ is slightly less than $2$ when $t$ is slightly bigger than $0$. 

    This means that the approximation is valid slightly below $x=2$.

    \qed
\end{solution}

\begin{example}
    A curve is defined by 
    \begin{align*}
        x &= \frac{1}{t} + 1, \; t\neq0\\
        y &= \cos t
    \end{align*}
    Show that we can write 
    \[ y = \cos\left(\frac{1}{ax + b}\right)\]

    Where $a$ and $b$ are constants to be found. Then find the derivative:
    \[ \frac{dy}{dx} \]
\end{example}

\begin{solution}
    Firstly we can rearrange $x$ in terms of $t$:
    \begin{align*}
        x &= \frac{1}{t} + 1 \\
        x -1 &= \frac{1}{t} \\
        t(x-1) &= 1 \\
        t &= \frac{1}{x-1} \tag{as long as $x\neq 1$}
    \end{align*}
    Then we can substitute $t$ into $y$ to find:
    \begin{align*}
     y &= \cos\left(\frac{1}{x-1} \right) 
    \end{align*}

    Which finds that $a=1$ and $b=-1$ so that 
    \[ y = \cos\left(\frac{1}{ax + b}\right)= \cos\left(\frac{1}{x-1}\right)\]

    Now to differentiate $y$ in terms of $x$, we must use the chain rule. We can ask our question of: \emph{if we could replace any part with $u$, how can we turn it into a function we know how to differentiate?}

    Note that we know how to differentiate $\cos u$, so we write:

    \begin{align*}
        y &= \cos\left(\frac{1}{x-1}\right) \\
        y &= \cos u \tag{what we want} \\
        u &= \frac{1}{x-1} \tag{to get what we want}
    \end{align*}

    Now remind ourselves of the chain rule Section \ref{sec:chain-rule}, e.g. by pretending that the $du$ cancel:
    \[ \frac{dy}{dx} = \frac{dy }{du}\frac{du }{dx}\]
    And therefore we can differentiate the equations to find the right-hand side of the above:
    \[ \frac{dy }{du} = -\sin u\]
    But what about differentiating $ u = \frac{1}{x-1}$? For this, we can do another chain rule! After asking ourselves the same question, we proceed with:

    \begin{align*}
        u &= \frac{1}{x-1}\\
        u &= \frac{1}{v} \tag{what we want}\\
        v &= x - 1 \tag{to get what we want}
    \end{align*}

    And the chain rule in this case would be:
    \[ \frac{du}{dx} = \frac{du }{dv}\frac{dv }{dx}\]

    We can differentiate each equation to get the right-hand side of the above:
    \begin{align*}
        \frac{du }{dv} &= -\frac{1}{v^2} \tag{writing $\frac{1}{v} = v^{-1}$}\\
        \frac{dv }{dx} &= 1 \\
        \therefore \frac{du}{dx} &= -\frac{1}{v^2} \tag{using the chain rule}\\
        &= -\frac{1 }{(x-1)^2} \tag{using $v = x-1$}
    \end{align*}
    Therefore, reminding ourselves, we have found:

    \begin{align*}
        \frac{du}{dx} &= -\frac{1 }{(x-1)^2} \\
        \frac{dy }{du} &= -\sin u \tag{from above}\\
        \therefore \frac{dy }{dx} &= -\sin u (-\frac{1 }{(x-1)^2}) \tag{using the chain rule}\\
        &= \frac{\sin u }{(x-1)^2} \\
        &= \frac{\sin \left( \frac{1}{x-1}\right)}{(x-1)^2} \tag{using $u =  \frac{1}{x-1}$}
    \end{align*}

    Therefore, if $y = \cos\left(\frac{1}{x-1}\right)$, then 

    \[\frac{dy }{dx} = \frac{\sin \left( \frac{1}{x-1}\right)}{(x-1)^2}\]

    This completes the question.

    \qed
\end{solution}


